%!TEX TS-program = lualatex
%!TEX encoding = UTF-8 Unicode

\documentclass[12pt, addpoints]{exam}

\printanswers

\usepackage[left=1in,right=0.75in,top=1in]{geometry}     

\usepackage{fontspec}
\def\mainfont{Linux Libertine O}
\setmainfont[Ligatures={TeX, Common}, BoldFont={* Bold}, ItalicFont={* Italic}, Numbers={Proportional, OldStyle}]{\mainfont}
%\setmonofont[Scale=MatchLowercase]{Inconsolata} 
\setsansfont[Scale=MatchLowercase]{Linux Biolinum O} 
\usepackage{microtype}

% This defines \amper for the fancy ampersand
% to be used in the header. See
% https://tex.stackexchange.com/a/58185/39194
\usepackage{xspace}
\newfontfamily\amperfont[Style=Alternate]{Linux Libertine O}
\makeatletter
\DeclareRobustCommand{\amper}{{\amperfont\ifx\f@shape\scname\smaller[1.2]\fi\&}\xspace}
\makeatother

\usepackage{array}
\newcolumntype{L}[1]{>{\raggedright\let\newline\\\arraybackslash\hspace{0pt}}p{#1}}
\newcolumntype{C}[1]{>{\centering\let\newline\\\arraybackslash\hspace{0pt}}p{#1}}
\newcolumntype{R}[1]{>{\raggedleft\let\newline\\\arraybackslash\hspace{0pt}}p{#1}}

\usepackage{booktabs}

\usepackage{graphicx}
	\graphicspath{%
	{/Users/goby/Pictures/teach/466/exams/}}%

% Used to get the last two digits of the year for the header below.
\def\short#1{\csname @gobbletwo\expandafter\endcsname\number#1}

\pagestyle{headandfoot}
\firstpageheader{Ornithology, Exam 1, \textsc{s}\short{\year}, \numpoints~points.}{}{%
	\ifprintanswers\textbf{KEY}\else Name: \enspace \makebox[2.5in]{\hrulefill}\fi}
\runningheader{}{}{\small{pg. \thepage}}
\runningheadrule

\footer{}{}{\iflastpage{\small \rule{0.6in}{0.4pt} / \numpoints\ points}{}}%{\makebox[0.75in]{\hrulefill}}
\extrafootheight{-0.5in}

\pointsinmargin
\marginpointname{ pts}

\usepackage{multicol}

%% Remove comment %% to print answer key.
\unframedsolutions
\renewcommand{\solutiontitle}{}
\SolutionEmphasis{\bfseries}

%% Create a Matching question format
\newcommand*\Matching[1]{
\ifprintanswers
	\textbf{#1}
\else
	\rule{2in}{0.5pt}
\fi
}
\newlength\matchlena
\newlength\matchlenb
\settowidth\matchlena{\rule{2.1in}{0pt}}
\newcommand\MatchQuestion[2]{%
	\setlength\matchlenb{0.98\linewidth}
	\addtolength\matchlenb{-\matchlena}
	\parbox[t]{\matchlena}{\Matching{#1}}\enspace\parbox[t]{\matchlenb}{#2}}


\newcommand*\AnswerBox[2]{%
    \parbox[t][#1]{0.92\textwidth}{%
    \begin{solution}#2\end{solution}}
    \vspace{\stretch{1}}
}

\newenvironment{AnswerPage}[1]
    {\begin{minipage}[t][#1]{0.92\textwidth}%
    \begin{solution}}
    {\end{solution}\end{minipage}
    \vspace{\stretch{1}}}

\newlength{\basespace}
\setlength{\basespace}{5\baselineskip}

\begin{document}

\begin{questions}


\vspace*{-1\baselineskip}

\question[5]
Match each letter to the correct structure. 1~point per blank. 

\vspace*{2ex}

\begin{multicols}{2}

%\textsc{a}.~\ifprintanswers\textbf{Wavelength}\else \rule{6cm}{0.4pt} \fi\\[1ex]

\textsc{a}.~\ifprintanswers\textbf{Rachis}\else \rule{5cm}{0.4pt} \fi 

\medskip%\\[1ex]

\textsc{b}.~\ifprintanswers\textbf{Distal barbule}\else \rule{5cm}{0.4pt} \fi %\\[1ex]

\medskip

\textsc{c}.~\ifprintanswers\textbf{Proximal barbule}\else \rule{5cm}{0.4pt} \fi %\\[1ex]

\medskip

\textsc{d}.~\ifprintanswers\textbf{Barbicel / hooklet}\else \rule{5cm}{0.4pt} \fi %\\[1ex]

\medskip

\textsc{e}.~\ifprintanswers\textbf{Barb}\else \rule{5cm}{0.4pt} \fi

\columnbreak
	\includegraphics[width=0.47\textwidth]{exam1_flight_vane}

\end{multicols}


\question[4]
Briefly describe the roles of the carina, the pectoralis, and the supracoracoideus during flight.

\AnswerBox{0.1\textheight}{Carina provides attachment point for pectoralis and supracoracoideus. Pectoralis is the large muscle that pulls the wings downward when it contracts. It is the power stroke. The smaller supracoracoideus pulls the wing upward when it contracts.}

%\question[5]
%Briefly list two skeletal structures of the skull that separate birds (and reptiles) from mammals. Tell how each differs from mammals. 
%
%\AnswerBox{0.1\textheight}{%
%One occipital condyle in birds vs 2 in mammals.\\
%One middle ear bone in birds vs 3 in mammals.\\
%Articular-quadrate jaw joint in birds vs dentary-squamosal in mammals.
%}%
	
	
\question[6]
Complete this table of comparisons between birds and mammals. Enter the number or name in the blanks provided. 1 point per blank.

\ifprintanswers
\begin{tabular}{L{0.75in}C{1.25in}C{1.25in}C{2.5in}}
	\toprule
	 & Number of occipital condyles & Number of middle ear bones & Type of jaw joint\tabularnewline
	\midrule
	& & &\\[0.35em]
	Birds & 1 & 1 & articular-quadrate \\[1.25em]
	Mammals & 2 & 3 & dentary-squamosal \tabularnewline
	\bottomrule
\end{tabular}
\else

\begin{tabular}{L{0.75in}C{1.25in}C{1.25in}C{2.5in}}
	\toprule
	 & Number of occipital condyles & Number of middle ear bones & Type of jaw joint\tabularnewline
	\midrule
	& & &\\[0.35em]
	Birds & \rule{1in}{0.4pt} & \rule{1in}{0.4pt}& \rule{2.4in}{0.4pt} \\[1.25em]
	Mammals & \rule{1in}{0.4pt} & \rule{1in}{0.4pt}& \rule{2.4in}{0.4pt} \tabularnewline
	\bottomrule
\end{tabular}

\fi

\question[5]
Briefly describe how counter-current exchange helps keep the body warm in birds that land regularly on cold surfaces.

\AnswerBox{0.1\textheight}{%
	Leg artery and vein closely pressed together. As warm blood from body passes down artery towards foot, heat diffuses to the colder blood in the vein. The heat is carried back to the body, keeping it warm. The arterial blood is always warmer than the adjacent venous blood.}

\newpage

\question[10]

\begin{multicols}{2}
Describe how birds change the angle of attack of the wing tips (lower arrows) and wing base (upper arrows) during the downstroke (A) and upstroke (B). Include how the angle of attack affects both thrust and lift for each stroke type.

\columnbreak

\includegraphics[width=\linewidth]{exam1_flight_upstroke_downstroke}
\end{multicols}


	\begin{AnswerPage}{0.3\textheight}%
	
	The angle of attack near the tips is such that lift and thrust are maximized on the downstroke. On the upstroke, the angle of attack at the tips is steeper, which reduces lift but also reduces reverse thrust. This angle reduces resistance as the bird raises its wing. 
	
	\bigskip
	
	At the base, the angle of attack does not change much and provides lift and thrust through most of the wing stroke.
		
	\end{AnswerPage}


\question[10]
Describe the complete respiratory cycle of birds. Include the air sacs and parabronchi in your description.

	\AnswerBox{0.3\textheight}{%
	Answer goes here.}

\newpage

%\question[10]
%Explain how having producer (finder) and scrounger individuals leads to an evolutionary stable strategy that maintains both foraging behaviors in the flock.
%	\AnswerBox{0.6\textheight}{Answer goes here.}%

\question[10]
Compare and contrast the arboreal hypothesis and the ontogenetic transition wing hypothesis proposed for the evolution of flight in birds.

	\AnswerBox{0.4\textheight}{Answer goes here.}%

%\question[5]
%Briefly describe competitive dominance structure in a foraging guild of tropical, ant-feeding birds. You don't have to name the types of birds, just describe the structure relative to the location (concentration) of ants.

%\question[5]
%Name the class and infraclass represented by the palate structure in the figure below. Describe the relationship of the vomer with either the maxillopalatine or the pterygoid to explain your answer.
%
%\hfill \includegraphics[width=0.4\textwidth]{palaeognath_jaw}


\question[10]
Name and describe the three types of soaring used by large birds to conserve energy during flight. (Not all soaring types are used by the same species of birds.)

\begin{AnswerPage}{0.3\textheight}
	Thermal soaring:
	
	Slope soaring:
	
	Dynamic soaring:
\end{AnswerPage}

%\question
%\begin{parts}
%\part[3] Name the two characteristics that define the Clade Avialae. Circle the one character that distinguishes modern birds from the non-avian theropod dinosaurs.
%
%\AnswerBox{0.1\textheight}{\textit{Powered flight} and feathers.}
%
%\part[4] Name the two infraclasses that make up the class Neornithes.
%
%\AnswerBox{0.1\textheight}{Neognathae and Paleognathae}
%
%\part[4] Name two of the three clades that make up the superorder Neoaves.
%
%\AnswerBox{0.1\textheight}{Passerea, Columbimorphae, and Mirandornithes.}
%
%\end{parts}

	
%\question[5]
%\begin{multicols}{2}
%Cory's Shearwater is an oceanic seabird that displays Lévy's Flight pattern when foraging. Briefly describe the flight pattern and explain how it increases the chance of foraging success for this species.

%\columnbreak

%\hfil \includegraphics[width=0.5\linewidth]{corys_shearwater} \hfill
%\end{multicols}


%\AnswerBox{0.1\textheight}{Answer goes here.}

\newpage

\question[15]
Below left is an Egyptian Vulture. Below right is a Magnificent Frigatebird. Assume both birds have the same wingspan length and the bodies have the same mass. Tell which bird has the highest aspect ratio and why (5 pts). Tell which bird has the highest wing loading and why (5 pts). Tell how aspect ratio and wing loading affect soaring ability and maneuverability (5 pts).

\includegraphics[width=\linewidth]{exam1_aspect_loading}

	\begin{AnswerPage}{0.4\textheight}%
	Answer goes here.

	\vspace*{\baselineskip}
	

	\end{AnswerPage}



\end{questions}


\end{document}